% ========================================
% ABSTRAKT 
% ========================================

\renewcommand{\rmdefault}{ptm} % Set font to Times New Roman
\renewcommand{\sfdefault}{phv} % Set sans-serif font to Helvetica
\chapter*{\fontsize{12}{14}\selectfont Abstrakt}

\vspace{0.5cm}
\noindent
\begin{tabular}{|p{5cm}|p{10cm}|}
\hline
\textbf{Imiona i nazwiska autorów projektu} & Franciszek Dębski, Patryk Wilk, Krystian Zając\\
\hline
\textbf{Rok urodzenia autorów projektu} & 2003/2004\\
\hline
\textbf{Stopień/tytuł zawodowy, imię i nazwisko prowadzącego przedmiotu} & mgr inż. Artur Pelo\\
\hline
\textbf{Instytut/Kierunek/Poziom/ Moduł specjalnościowy} & Instytut Nauk Inżynieryjno-Technicznych / Informatyka / Inżynierskie / Informatyka w przedsiębiorstwie\\
\hline
\textbf{Data realizacji projektu} & od 3.03.2026 do 9.06.2026\\
\hline
\textbf{Nazwa przedmiotu} & projekt zespołowy\\
\hline
\textbf{Rok studiów/semestr} & rok III, semestr VI\\
\hline
\end{tabular}

\vspace{1.5cm}

\noindent
\begin{tabular}{|p{5cm}|p{10cm}|}
\hline
\textbf{Tytuł projektu} & Aplikacja w architekturze mikroserwisowej\\
\hline
\textbf{Słowa kluczowe (max 5)} & informatyka, aplikacje mobilne, bazy danych\\
\hline
\textbf{Streszczenie (max 1400 znaków)} & Aplikacja w architekturze mikroserwisowej umożliwiająca integrację usług i komponentów w celu tworzenia skalowalnych i utrzymywanych systemów komunikacji między IoT a aplikacją webową i mobilną. Projekt obejmuje implementację mikroserwisów (serwer, baza danych), interfejsu użytkownika(Web) oraz integrację z urządzeniami IoT(ESP32) z czujnikiem. \\
\hline
\end{tabular}

\vspace{1.5cm}
